% =============================================================================
% Vehicle Parameter Identification — Report Outline
% =============================================================================
% Can be compiled standalone or included in your main report document.
% Data CSV files are saved by the test scripts in f1tenth_parameters/data/.

\documentclass[11pt,a4paper]{article}
\usepackage{amsmath,amssymb}
\usepackage{siunitx}
\usepackage{booktabs}
\usepackage{graphicx}
\graphicspath{{../figures/}}
\usepackage[margin=2.5cm]{geometry}
\usepackage{hyperref}

% Conditionally include a figure only if the PDF exists
\newcommand{\includefigure}[3][width=0.7\textwidth]{%
  \IfFileExists{../figures/#2}{%
    \begin{figure}[htbp]
      \centering
      \includegraphics[#1]{#2}
      \caption{#3}
    \end{figure}%
  }{\par\textbf{[Figure not yet generated: \detokenize{#2}]}\par}%
}

\title{Vehicle Parameter Identification --- Report Outline}
\author{}
\date{}

\begin{document}
\maketitle

\section{Vehicle Parameter Identification}\label{sec:param-id}

The dynamic bicycle model used by the MPC controller requires physical
parameters that cannot be taken from datasheets alone.
This section describes the experimental procedures used to identify
them on the real F1/10th vehicle.

All tests use a common ROS\,2 framework that records time-stamped
sensor data at approximately \SI{200}{\hertz}:
VESC odometry (speed from ERPM), VESC IMU (3-axis accelerometer
and gyroscope), and motor state (current, voltage, temperature).
Each test saves raw data to CSV for post-processing and automatically
updates a central \texttt{vehicle\_params.yaml} file.

%%% ===================================================================
\subsection{Dynamic Bicycle Model}\label{sec:bicycle-model}

The MPC uses a 7-state dynamic bicycle model with state
$\mathbf{x} = [x,\, y,\, \psi,\, v_x,\, v_y,\, \omega,\, \omega_w]$
and control inputs $\mathbf{u} = [\delta,\, T_\text{motor}]$.
% TODO: brief sentence about why the dynamic (not kinematic) model is needed.
The equations of motion are:
\begin{align}
  \dot{v}_x &= \frac{1}{m}\bigl(F_x - F_{yf}\sin\delta + m\,v_y\,\omega\bigr)
  \label{eq:dvx}\\
  \dot{v}_y &= \frac{1}{m}\bigl(F_{yf}\cos\delta + F_{yr} - m\,v_x\,\omega\bigr)
  \label{eq:dvy}\\
  \dot{\omega} &= \frac{1}{I_z}\bigl(l_f\,F_{yf}\cos\delta - l_r\,F_{yr}\bigr)
  \label{eq:domega}\\
  \dot{\omega}_w &= \frac{1}{I_w}\Bigl(\frac{T_\text{motor}}{G} - F_x\,R_w\Bigr)
  \label{eq:domegaw}
\end{align}
where $G$ is the gear ratio, $R_w$ the wheel radius, and $I_w$ the
drivetrain inertia. Lateral tire forces use a linear model:
\begin{equation}\label{eq:tire-lat}
  F_{yf} = C_{\alpha f}\,\alpha_f, \qquad
  F_{yr} = C_{\alpha r}\,\alpha_r
\end{equation}
with slip angles
\begin{equation}\label{eq:slip-angles}
  \alpha_f = \delta - \arctan\!\frac{v_y + l_f\,\omega}{v_x}, \qquad
  \alpha_r = -\arctan\!\frac{v_y - l_r\,\omega}{v_x}.
\end{equation}
The longitudinal tire force is $F_x = C_x\,\kappa$, where $\kappa$ is the
slip ratio. Table~\ref{tab:params} lists all model parameters and the test
used to identify each one.

%%% ===================================================================
\subsection{Tire Friction Coefficient}\label{sec:friction}

\paragraph{Theory.}
At steady-state circular driving the centripetal acceleration is
$a_y = v^2/R$. The maximum sustainable $a_y$ before the tires begin to
slide gives the friction coefficient:
\begin{equation}\label{eq:mu}
  \mu = \frac{a_{y,\max}}{g}.
\end{equation}
The IMU measures $a_y$ directly, independent of odometry errors.

\paragraph{Procedure.}
The car drives circles at a fixed steering angle with speed increasing
in \SI{0.5}{\metre\per\second} steps from \SIrange{1.0}{4.0}{\metre\per\second}.
At each speed, \SI{8}{\second} of steady-state IMU data is recorded.

\paragraph{Data.}
\includefigure{friction_ay_vs_speed.pdf}{%
  Lateral acceleration (IMU) vs.\ speed.  The friction coefficient
  $\mu$ is determined from the plateau.\label{fig:friction}}
\textbf{Table}: speed $\mid$ $\bar{a}_y \pm \sigma$ $\mid$ $\mu=a_y/g$.

\noindent CSV columns used: \texttt{imu\_ay}, \texttt{speed\_setpoint}, \texttt{phase}.
Group by \texttt{speed\_setpoint}, compute mean and std of \texttt{imu\_ay}
in the \texttt{recording} phase.
Before aggregation, use bias-corrected IMU channels and keep only
near steady-state samples where $|a_y|\approx|v_x\,\omega|$.

\paragraph{Uncertainty.}
\begin{itemize}
  \item IMU bias: \SI{2}{\second} stationary calibration; residual
        $< \SI{0.05}{\metre\per\second\squared}$.
  \item Surface-dependent: result only valid for the tested floor.
  \item Temperature-dependent: rubber grip varies with tire temperature.
\end{itemize}

%%% ===================================================================
\subsection{Cornering Stiffness}\label{sec:cornering}

\paragraph{Theory.}
In steady-state cornering $(\dot{v}_y=\dot{\omega}=0)$, the front and
rear lateral forces and slip angles reduce to:
\begin{equation}\label{eq:Fy-ss}
  F_{yf} = m\,a_y\,\frac{l_r}{L}, \qquad
  F_{yr} = m\,a_y\,\frac{l_f}{L},
\end{equation}
\begin{equation}\label{eq:alpha-ss}
  \alpha_f \approx \delta - \frac{l_f\,\omega}{v_x}, \qquad
  \alpha_r \approx \frac{l_r\,\omega}{v_x},
\end{equation}
where $\omega$ is the IMU yaw rate and $v_x$ the ERPM-based speed.
The cornering stiffnesses follow from
$C_{\alpha f} = F_{yf}/\alpha_f$ and $C_{\alpha r} = F_{yr}/\alpha_r$.
The understeer gradient is
\begin{equation}\label{eq:Kus}
  K_\text{us} = \frac{m\,g}{L}\!\left(\frac{l_f}{C_{\alpha f}}
  - \frac{l_r}{C_{\alpha r}}\right).
\end{equation}

\paragraph{Procedure.}
Same steady-state circles as the friction test, at speeds from
\SIrange{1.0}{3.5}{\metre\per\second}. At each speed, \SI{5}{\second} of
$v_x$, $\omega$ (IMU), and $a_y$ (IMU) are recorded.

\paragraph{Data.}
\includefigure[width=0.48\textwidth]{cornering_Fy_vs_alpha.pdf}{%
  Lateral force vs.\ slip angle for front and rear axles.  The linear
  region yields the cornering stiffness $C_\alpha$.\label{fig:cornering-Fy}}

\includefigure[width=0.48\textwidth]{cornering_alpha_vs_speed.pdf}{%
  Slip angle growth with speed for front and rear axles.\label{fig:cornering-alpha}}
\textbf{Table}: speed $\mid$ $\alpha_f$ $\mid$ $\alpha_r$ $\mid$
$C_{\alpha f}$ $\mid$ $C_{\alpha r}$.

\noindent CSV columns used: \texttt{odom\_vx}, \texttt{imu\_gz},
\texttt{imu\_ay}, \texttt{cmd\_steering}, \texttt{phase}.
Compute $\alpha_f$, $\alpha_r$ from \eqref{eq:alpha-ss}, then
$F_{yf}$, $F_{yr}$ from \eqref{eq:Fy-ss}.
Bias-corrected IMU data is used, and samples with large
$|a_y-v_x\omega|$ consistency error are filtered out.

\paragraph{Uncertainty.}
\begin{itemize}
  \item Relies on measured $m$, $l_f$, $l_r$ ($\pm 1$--$\SI{2}{\milli\metre}$).
  \item Small-angle approximation breaks down above $\sim 5^\circ$ slip.
  \item IMU yaw-rate bias drift (mitigated by \SI{2}{\second} calibration).
  \item ERPM noise affects $v_x$ and thus slip-angle calculation.
\end{itemize}

%%% ===================================================================
\subsection{Longitudinal Tire Stiffness}\label{sec:long-stiffness}

\paragraph{Theory.}
The slip ratio is
\begin{equation}\label{eq:kappa}
  \kappa = \frac{v_\text{wheel} - v_\text{body}}
               {\max(|v_\text{wheel}|,\,|v_\text{body}|)},
\end{equation}
where $v_\text{wheel}$ comes from ERPM and $v_\text{body}$ from an
independent velocity sensor. In the linear region ($|\kappa|<0.15$):
\begin{equation}\label{eq:Fx}
  F_x = C_x\,\kappa, \qquad F_x = m\,a_x.
\end{equation}

\paragraph{Limitation.}
Because the only available body-velocity estimate is obtained by
integrating the IMU, and the accelerometer is affected by pitch
oscillations during braking, $v_\text{body}$ drifts rapidly.
As a result, the slip ratio $\kappa$ is dominated by integration error and
$C_x$ cannot be reliably identified from the current data.
A lidar-based scan-matching odometry would provide an independent
body velocity to resolve this issue.

\paragraph{Traction force balance.}
Instead, the motor torque test data are used to verify the
force-transmission chain.  The wheel-side force is computed from motor
current as
$F_\text{motor} = K_t\,G / r_\text{eff}\cdot I$,
while the body-side force is $F_\text{body} = m\,a_x$ from the IMU.
In the unsaturated region, the linear slope gives the drivetrain
efficiency $\eta$.  At high motor current the body acceleration
saturates, revealing the longitudinal traction limit.

\paragraph{Data.}
\includefigure[width=0.48\textwidth]{long_v_comparison.pdf}{%
  $v_\text{wheel}$ (ERPM) and $v_\text{body}$ (IMU) vs.\ time.
  The gap between them is the longitudinal slip.\label{fig:long-v}}

\includefigure[width=0.48\textwidth]{long_Fx_vs_kappa.pdf}{%
  Motor force vs.\ body force from the motor torque test.
  The linear slope gives the drivetrain efficiency; the saturation
  plateau indicates the traction limit $\mu_x$.\label{fig:long-Fx}}

\noindent CSV sources: \texttt{longitudinal\_stiffness} (velocity comparison),
\texttt{motor\_torque} (force balance).

\paragraph{Uncertainty.}
\begin{itemize}
  \item IMU integration drift accumulates; slip ratio data is unreliable.
  \item 4WD distributes torque through differentials; efficiency
        $\eta$ includes differential and drivetrain losses.
  \item Traction limit depends on surface and tire temperature.
\end{itemize}

%%% ===================================================================
\subsection{Motor Torque Mapping}\label{sec:motor-torque}

\paragraph{Theory.}
Two independent estimates of wheel force:
\begin{align}
  F_\text{current} &= K_t\,I_\text{motor}\,\frac{G}{r_\text{eff}},
  \label{eq:F-current} \\
  F_\text{accel}   &= m\,a_x \quad\text{(IMU)},
  \label{eq:F-accel}
\end{align}
where $K_t = 60/(2\pi\,K_v)$ and $G=11.82$ is the gear ratio.
Plotting $F_\text{accel}$ against $I_\text{motor}$ yields the effective
current-to-force gain (slope) and drivetrain losses (intercept).

\paragraph{Procedure.}
Accelerate to target speeds (\num{1.5}, \num{2.0}, \num{3.0},
\SI{4.0}{\metre\per\second}), recording motor current and IMU $a_x$
simultaneously. Braking (regenerative) current is measured when the car
decelerates.

\paragraph{Data.}
\includefigure[width=0.48\textwidth]{torque_I_vs_F.pdf}{%
  Motor current vs.\ net force with linear fit.\label{fig:torque-IF}}

\includefigure[width=0.48\textwidth]{torque_current_vs_time.pdf}{%
  Motor current time series at each speed target.\label{fig:torque-time}}
\textbf{Table}: speed $\mid$ $\bar{I}$ $\mid$ peak $I$ $\mid$
$F_\text{net}$ $\mid$ efficiency.

\noindent CSV columns used: \texttt{motor\_current}, \texttt{imu\_ax},
\texttt{cmd\_speed}, \texttt{phase}.
Acceleration is computed from bias-corrected IMU $a_x$.

\paragraph{Uncertainty.}
\begin{itemize}
  \item Motor $K_v$ from datasheet ($\pm5\%$ typical for hobby motors).
  \item $r_\text{eff}$ measurement accuracy ($\pm\SI{1}{\milli\metre}$
        $\Rightarrow \sim2\%$ error on force).
  \item Gear efficiency not measured separately; absorbed into effective $K_t$.
\end{itemize}

%%% ===================================================================
\subsection{Maximum Dynamics}\label{sec:max-dynamics}

\paragraph{Theory.}
Direct measurement of the vehicle's performance envelope:
$v_\text{max}$, $a_\text{max}$ (peak acceleration),
and $a_\text{brake}$ (peak deceleration).
These define the MPC input/state constraints.

\paragraph{Procedure.}
Command maximum speed for \SI{5}{\second} (acceleration phase),
then command zero (braking phase). Both ERPM-based speed and IMU $a_x$
are recorded. A velocity plateau detector checks whether the car
actually reached top speed.

\paragraph{Data.}
\includefigure{max_dyn_speed_vs_time.pdf}{%
  Speed (odometry, blue) and IMU-integrated speed (orange) vs.\ time.
  The braking region shows odometry dropping to zero prematurely due to
  wheel slip.\label{fig:maxdyn}}

\noindent CSV columns used: \texttt{odom\_vx}, \texttt{imu\_ax},
\texttt{phase}.
Both peak acceleration/deceleration and IMU-integrated speed use
bias-corrected IMU $a_x$.

\paragraph{Uncertainty.}
\begin{itemize}
  \item Odom speed unreliable during braking (wheel lock/slip); use IMU.
  \item IMU-integrated velocity drifts; odom more reliable during steady
        acceleration.
  \item Battery voltage sag under load reduces effective $v_\text{max}$.
\end{itemize}

%%% ===================================================================
\subsection{Steering Actuator Rate}\label{sec:steering-rate}

\paragraph{Theory.}
The servo has a finite angular velocity. By commanding step changes
$0 \to \delta_\text{max}$ and observing the 10\%--90\% rise time of
the IMU yaw rate:
\begin{equation}\label{eq:steer-rate}
  \dot{\delta}_\text{max} = \frac{\Delta\delta}{\Delta t_{10-90}}.
\end{equation}

\paragraph{Procedure.}
Drive at constant \SI{1.5}{\metre\per\second}. Alternate steering
between $+\delta$ and $-\delta$ for 6 transitions. Record IMU yaw rate
$\omega_z$ and commanded steering angle.

\paragraph{Data.}
\includefigure{steering_step_response.pdf}{%
  IMU yaw rate (solid) and commanded steering (dashed) vs.\ time.
  The slope during transitions gives the effective steering
  rate.\label{fig:steer}}
\textbf{Table}: transition $\mid$ $\Delta\delta$ $\mid$ $\Delta t$ $\mid$
rate (\si{\radian\per\second}).

\noindent CSV columns used: \texttt{imu\_gz}, \texttt{cmd\_steering},
\texttt{phase}.

\paragraph{Uncertainty.}
\begin{itemize}
  \item Yaw rate includes vehicle dynamics response, not only servo motion.
  \item At low speed ($\SI{1.5}{\metre\per\second}$) the dynamic lag is
        small relative to the servo speed.
  \item Servo speed may depend on battery voltage.
\end{itemize}

%%% ===================================================================
\subsection{Parameter Summary}\label{sec:param-summary}

Table~\ref{tab:params} compares tested values with the F1/10th gym
simulation defaults.

\begin{table}[h]
\centering
\caption{Identified vehicle model parameters.}\label{tab:params}
\begin{tabular}{ll S[table-format=2.3] S[table-format=2.3] l}
\toprule
{Parameter} & {Symbol} & {Measured} & {Sim default} & {Unit} \\
\midrule
Friction coeff.      & $\mu$                    & {---} & 1.049  & {--}             \\
Front corn.\ stiff.  & $C_{\alpha f}$           & {---} & 4.718  & \si{\newton\per\radian} \\
Rear corn.\ stiff.   & $C_{\alpha r}$           & {---} & 5.456  & \si{\newton\per\radian} \\
Long.\ stiffness     & $C_x$                    & N/A\textsuperscript{*} & 300    & \si{\newton}     \\
Max drive torque     & $T_\text{max}$            & {---} & 22.9   & \si{\newton\metre}\\
Max brake torque     & $T_\text{min}$            & {---} & {-24.1}& \si{\newton\metre}\\
Max velocity         & $v_\text{max}$            & {---} & 20.0   & \si{\metre\per\second} \\
Max acceleration     & $a_\text{max}$            & {---} & 9.51   & \si{\metre\per\second\squared} \\
Max deceleration     & $a_\text{brake}$          & {---} & 10.0   & \si{\metre\per\second\squared} \\
Max steer.\ rate     & $\dot{\delta}_\text{max}$ & {---} & 3.2    & \si{\radian\per\second} \\
\bottomrule
\end{tabular}

\smallskip\noindent\textsuperscript{*}$C_x$ could not be identified:
IMU-integrated body velocity drifts too much for reliable slip-ratio
estimation.  Lidar-based scan-matching odometry is recommended.
\end{table}

% TODO: Fill in the "Measured" column after running the tests.
% TODO: Discuss which parameters differ significantly from the sim defaults
%       and what implications this has for MPC tuning.
% TODO: Note which sensor (IMU vs odometry) proved more reliable for each test.
% TODO: Note surface dependency: friction and stiffness change with floor type.

\end{document}
